\section{产生湮灭算符和基本性质}

\subsection{定义与代数}

设 $\{\lvert\alpha\rangle\}$ 为单粒子完备正交归一基。玻色子产生、湮灭算符满足
\begin{equation}
  [b_\alpha,b_\beta^\dagger]=\delta_{\alpha\beta},
  \qquad
  [b_\alpha,b_\beta]=[b_\alpha^\dagger,b_\beta^\dagger]=0;
\end{equation}
费米子满足反对易关系
\begin{equation}
  \{c_\alpha,c_\beta^\dagger\}=\delta_{\alpha\beta},
  \qquad
  \{c_\alpha,c_\beta\}=\{c_\alpha^\dagger,c_\beta^\dagger\}=0.
\end{equation}
粒子数算符 $n_\alpha=a_\alpha^\dagger a_\alpha$（$a$ 代表 $b$ 或 $c$）的本征值为占据数。

对费米子，若约定模式按固定次序排列，则
\begin{equation}
  c_\alpha\lvert\{n\}\rangle
  =(-1)^{S_\alpha}\sqrt{n_\alpha}\,\lvert n_\alpha-1\rangle,
  \qquad
  c_\alpha^\dagger\lvert\{n\}\rangle
  =(-1)^{S_\alpha}\sqrt{1-n_\alpha}\,\lvert n_\alpha+1\rangle,
\end{equation}
其中 $S_\alpha=\sum_{\beta<\alpha}n_\beta$ 是相位因子，用以保持整体反对称。

\subsection{场算符与密度}

在坐标表象中，
\begin{equation}
  \psi_\sigma(\mathbf{r})
  =\sum_\alpha \phi_\alpha(\mathbf{r})\,c_{\alpha\sigma},
  \qquad
  \psi_\sigma^\dagger(\mathbf{r})
  =\sum_\alpha \phi_\alpha^*(\mathbf{r})\,c_{\alpha\sigma}^\dagger,
\end{equation}
并满足
\begin{equation}
  \{\psi_\sigma(\mathbf{r}),\psi_{\sigma'}^\dagger(\mathbf{r}')\}
  =\delta_{\sigma\sigma'}\delta(\mathbf{r}-\mathbf{r}').
\end{equation}
粒子密度与粒子数算符为
\begin{equation}
  \hat n(\mathbf{r})=\sum_\sigma\psi_\sigma^\dagger(\mathbf{r})\psi_\sigma(\mathbf{r}),
  \qquad
  \hat N=\int\mathrm{d}\mathbf{r}\,\hat n(\mathbf{r})
  =\sum_{\alpha\sigma}c_{\alpha\sigma}^\dagger c_{\alpha\sigma}.
\end{equation}
平面波基 $\phi_{\mathbf{k}}(\mathbf{r})=V^{-1/2}e^{\mathrm{i}\mathbf{k}\cdot\mathbf{r}}$ 下，
\begin{equation}
  \psi_\sigma(\mathbf{r})
  =\frac{1}{\sqrt{V}}\sum_{\mathbf{k}}
  e^{\mathrm{i}\mathbf{k}\cdot\mathbf{r}}c_{\mathbf{k}\sigma},
\end{equation}
密度起伏的 Fourier 分量（Giuliani--Vignale 常用记号）为
\begin{equation}
  \hat n_{\mathbf{q}}
  =\sum_{\mathbf{k}\sigma}
  c_{\mathbf{k}\sigma}^\dagger c_{\mathbf{k}+\mathbf{q},\sigma}
  =\int\mathrm{d}\mathbf{r}\,e^{-\mathrm{i}\mathbf{q}\cdot\mathbf{r}}\hat n(\mathbf{r}).
\end{equation}

\subsection{单体算符的二次量子化：推导}

一阶量子化中单体算符 $A=\sum_{i=1}^N a(i)$。在 Fock 空间中，其对任意占据组态的矩阵元可写为
\begin{equation}
  A=\sum_{\alpha\beta}
  \langle\alpha|a|\beta\rangle\,c_\alpha^\dagger c_\beta.
\end{equation}
\textbf{推导要点}：在 $N$ 粒子反对称态上，把 $a(1)$ 夹在两 Slater 行列式之间，用 Laplace 展开与轨道正交性，只留下单粒子矩阵元 $\langle\alpha|a|\beta\rangle$；再用产生湮灭算符重写占据数的升降，即得上式。动能与外势因此成为
\begin{equation}
  T=\sum_{\mathbf{k}\sigma}\frac{\hbar^2k^2}{2m}
  c_{\mathbf{k}\sigma}^\dagger c_{\mathbf{k}\sigma},
  \qquad
  V_{\mathrm{ext}}
  =\sum_{\mathbf{k}\mathbf{k}'\sigma}
  \tilde v_{\mathrm{ext}}(\mathbf{k}-\mathbf{k}')
  c_{\mathbf{k}\sigma}^\dagger c_{\mathbf{k}'\sigma}.
\end{equation}

\subsection{二体相互作用：推导与矩阵元}

二体算符 $V=\frac12\sum_{i\neq j}v(\mathbf{r}_i-\mathbf{r}_j)$ 的二次量子化形式为
\begin{equation}
  V=\frac12\sum_{\alpha\beta\gamma\delta}
  \langle\alpha\beta|v|\gamma\delta\rangle
  c_\alpha^\dagger c_\beta^\dagger c_\delta c_\gamma,
\label{eq:V2nd}
\end{equation}
其中
\begin{equation}
  \langle\alpha\beta|v|\gamma\delta\rangle
  =\int\mathrm{d}\mathbf{r}\,\mathrm{d}\mathbf{r}'\,
  \phi_\alpha^*(\mathbf{r})\phi_\beta^*(\mathbf{r}')
  v(\mathbf{r}-\mathbf{r}')
  \phi_\gamma(\mathbf{r})\phi_\delta(\mathbf{r}').
\end{equation}
算符顺序取 $c^\dagger c^\dagger c\,c$（先湮灭后产生的反向书写），以使费米子符号与行列式展开一致。

对平移不变的库仑作用 $v(\mathbf{r})=e^2/r$，平面波矩阵元只在动量守恒时非零。令
\begin{equation}
  v_{\mathbf{q}}=\int\mathrm{d}\mathbf{r}\,e^{-\mathrm{i}\mathbf{q}\cdot\mathbf{r}}\frac{e^2}{r}
  =\frac{4\pi e^2}{q^2}
  \qquad(q\neq0),\label{eq:vq}
\end{equation}
则（自旋守恒、忽略 $q=0$ 项时）
\begin{equation}
  V=\frac{1}{2V}\sum_{\mathbf{q}\neq0}v_{\mathbf{q}}
  \sum_{\substack{\mathbf{k}\mathbf{k}'\\ \sigma\sigma'}}
  c_{\mathbf{k}+\mathbf{q},\sigma}^\dagger
  c_{\mathbf{k}'-\mathbf{q},\sigma'}^\dagger
  c_{\mathbf{k}',\sigma'}
  c_{\mathbf{k},\sigma}.
\label{eq:Vcoulomb}
\end{equation}
等价地，
\begin{equation}
  V=\frac{1}{2V}\sum_{\mathbf{q}\neq0}v_{\mathbf{q}}
  \bigl(\hat n_{\mathbf{q}}\hat n_{-\mathbf{q}}-\hat N\bigr),
\end{equation}
其中 $-\hat N$ 来自 $i=j$ 自作用的扣除（因 $\hat n_{\mathbf{q}}\hat n_{-\mathbf{q}}$ 含自作用项）。

\begin{example}[自由费米气]
无相互作用时
\begin{equation}
  H_0=\sum_{\mathbf{k}\sigma}\varepsilon_{\mathbf{k}}
  c_{\mathbf{k}\sigma}^\dagger c_{\mathbf{k}\sigma},
  \qquad
  \varepsilon_{\mathbf{k}}=\frac{\hbar^2k^2}{2m}.
\end{equation}
零温基态为费米海：$n_{\mathbf{k}\sigma}=\theta(k_F-k)$，密度
\begin{equation}
  n=\frac{N}{V}=\frac{k_F^3}{3\pi^2}
  \quad\Rightarrow\quad
  k_F=(3\pi^2 n)^{1/3}.
\end{equation}
\end{example}

\subsection{运动方程}

海森堡运动方程
\begin{equation}
  i\hbar\partial_t c_{\mathbf{k}\sigma}(t)
  =[c_{\mathbf{k}\sigma}(t),H]
\end{equation}
在含 \eqref{eq:Vcoulomb} 时，右边产生三个算符的乘积，从而关联函数满足 Martin--Schwinger 无穷层级。平均场截断（用密度平均值替换算符）即导向 Hartree--Fock。
